% Renamed from long space-filled RGE title
% Note: digit-level m_H claims calibrated to 125 GeV class in hardening layer.
\documentclass[11pt,a4paper]{article}
\usepackage[utf8]{inputenc}
\usepackage{amsmath,amssymb,amsfonts}
\usepackage{booktabs}
\usepackage{hyperref}
\usepackage{geometry}
\geometry{margin=2.5cm}

\title{\textbf{Two-Loop RG Analysis and Uncertainty Quantification in SAM3} \\ (v4.26 baseline; calibrated August 2026)}

\author{Shawn Dykes \\ SAM3 Collaboration}
\date{May 2026 / August 2026}

\begin{document}
\maketitle

\begin{center}
\fbox{\parbox{0.92\linewidth}{\textbf{Supersession notice.}
Digit-level $m_H=125.1\pm 0.4$\,GeV language is superseded by the hardening layer:
Higgs mass in the 125\,GeV class with a wider theoretical band under present residuals.}}
\end{center}

\vspace{1em}

\begin{abstract}
We summarize two-loop renormalization-group analysis of the SAM3 framework with partial higher-order corrections and uncertainty propagation. $\omega_0$ is geometrically derived. The Higgs mass is reported in the 125\,GeV class consistent with the August 2026 residual map.
\end{abstract}

\section{Setup}
Two-loop beta functions follow standard SM literature. High-scale boundary conditions are taken from the Dual-Zero spectral action at the geometric scale set by $\ell_0$, with $\omega_0\approx 0.927$.

\section{Results (calibrated)}
Running produces a Higgs mass in the 125\,GeV class. Monte-Carlo propagation of overlap, grid, and higher-order uncertainties supports a theoretical band of order a few GeV under present residuals; sub-GeV digit matching is not claimed.

\section{Conclusion}
RG analysis remains part of the predictive pipeline. Claim strength follows \texttt{docs/hardening/13\_Seeley\_DeWitt\_GN\_Higgs.md} and the repository status file.

\textbf{Repository}: \url{https://github.com/mohawksd9sd-maker/SAM3-DualZero-Conoid}

\end{document}
